\documentclass[11pt]{article}
\usepackage[utf8]{inputenc}
\usepackage[T1]{fontenc}
\usepackage{lmodern}
\usepackage{amsmath,amssymb,mathtools}
\usepackage[margin=1in]{geometry}
\usepackage{microtype}
\usepackage{enumitem}

\allowdisplaybreaks
\setlength{\parindent}{1.5em}
\setlength{\parskip}{0pt}

\newcommand{\cosec}{\operatorname{cosec}}
\newcommand{\Min}{\operatorname{Min}}
\newcommand{\Max}{\operatorname{Max}}
\newcommand{\dd}{\,d}

\title{\textit{The Lattice Points of a Circle}}
\author{J. E. Littlewood, F.R.S., and A. Walfisz\\
\large (With a Note by Prof. E. Landau.)}
\date{(Received April 22, 1924.)}

\begin{document}
\maketitle

\noindent\textbf{1.}
Let $r(x)$ denote the number of ways in which the positive integer $x$ can
be expressed as the sum of two squares (positive, negative or zero), and let
\[
 \mathrm{R}(x)=\sum_{0\leq n\leq x}r(n)
 =\sum_{0\leq p^2+q^2\leq x}1.
\]
Thus $\mathrm{R}(x)$ is the number of ``lattice-points'' (points whose
co-ordinates $p,q$ are integers, positive, negative or zero) in or on the
boundary of the circle with centre at the origin and radius $\sqrt{x}$. It is
trivial that
\begin{equation}
 \mathrm{R}(x)-\pi x=O\!\left(x^{1/2}\right),
 \tag{1.1}
\end{equation}
it has been shown by Hardy\footnote{G. H. Hardy, ``On the expression of a
number as the sum of two squares,'' \textit{Quarterly Journal of Math.},
vol.~46, pp.~263--283 (1915).} and Landau\footnote{E. Landau, ``Über die
Gitterpunkte in einem Kreise (II),'' \textit{Göttinger Nachrichten},
pp.~161--171 (1915).} that the relation
\[
 \mathrm{R}(x)-\pi x=O\!\left(x^\alpha\right)
\]
is not true for any constant $\alpha<\tfrac14$, and it has been known for
some time that
\[
 \mathrm{R}(x)-\pi x=O\!\left(x^{1/3}\right).
\]
A very important advance is due to van der Corput, who proved in a recent
paper\footnote{J. G. van der Corput, ``Neue zahlentheoretische
Abschätzungen,'' \textit{Math. Annalen}, vol.~89, pp.~215--254 (1923).} that
\begin{equation}
 \mathrm{R}(x)-\pi x=O\!\left(x^{\Theta+\varepsilon}\right),
 \tag{1.2}
\end{equation}
where $\Theta$ is some constant less than $\tfrac13$. For the corresponding
problem concerning
\[
 \mathrm{D}(x)=\sum_{0\leq n\leq x}d(n),
\]
where $d(n)$ is the number of divisors of $n$, van der Corput obtains\footnote{
J. G. van der Corput, ``Verschärfung der Abschätzung beim Teilerproblem,''
\textit{Math. Annalen}, vol.~87, pp.~39--65 (1922). A number of misprints in
this paper are corrected by the author in the errata attached to vol.~89.}
the more precise result
\begin{equation}
 \mathrm{D}(x)-\mathrm{D}_0(x)=O\!\left(x^{\Theta+\varepsilon}\right),
 \tag{1.3}
\end{equation}
where $\mathrm{D}_0(x)$ is a certain elementary function whose leading term
is $x\log x$, $\varepsilon$ is an arbitrary positive number, and $\Theta$ is
some constant less than $\tfrac{33}{100}$.

The present writers, applying a method used by Hardy and
Littlewood\footnote{See E. Landau, ``Über die $\zeta$-Funktion und die
$L$-Funktionen,'' \textit{Math. Zeitschr.}, vol.~20, pp.~105--125 (1924).}
to prove that
\begin{equation}
 \zeta\!\left(\tfrac12+it\right)=O\!\left(t^{1/6+\varepsilon}\right),
 \tag{1.4}
\end{equation}
discovered, independently, that it is possible to prove (1.2) with
\begin{equation}
 \Theta=\frac{37}{112}.
 \tag{1.5}
\end{equation}
The result (1.5) is new, but we do not regard this as the main point of the
present paper. Doubtless no proof of (1.2) (with $\Theta<\tfrac13$) can be
altogether easy; but whereas van der Corput's method is probably the most
formidable argument in the whole of pure mathematics, it will be found that
the proof given here is both reasonably short and not impossible to grasp in
its entirety.

Our proof rests on the possibility of finding an expression as a (finite or
infinite) series for $\mathrm{R}(x)-\pi x$, sufficiently approximate in itself,
and amenable (in essence) to the argument of \S\S~5--7 below. There are, in
point of fact, alternatives to the expression we adopt (which is based on a
result due to Wigert); the several versions of the proof all lead, however,
to the same value of $\Theta$.

\medskip
\noindent\textbf{2.}
We suppose always $x\geq2$, and write for convenience
\[
 s=x^{75/112}.
\]
The symbols $c_1,c_2,\ldots,c_{30}$ denote positive constants depending only
on $\alpha$, a positive constant introduced in \S~5, and $O$'s depend only on
$\varepsilon$, and on $\alpha$ when that is a parameter.

\medskip
\noindent\textsc{Lemma 1.} \textit{We have}
\[
 \mathrm{R}(x)-\pi x
 =\frac1\pi x^{1/4}\sum_{n=1}^{\infty}\frac{r(n)}{n^{3/4}}a_n
  +O\!\left(x^{37/112+\varepsilon}\right),
\]
\textit{where}
\begin{equation}
 a_n=\int_0^\infty e^{-u}u^{1/(4s)}
 \cos\!\left(2\pi\sqrt{nx}\,u^{1/(2s)}-\frac{3\pi}{4}\right)\dd u.
 \tag{2.1}
\end{equation}
This result is due to Wigert,\footnote{S. Wigert, ``Über das Problem der
Gitterpunkte in einem Kreise,'' \textit{Math. Zeitschr.}, vol.~5,
pp.~310--318 (1919).} who proves it, indeed, uniformly in $s\geq1$ and with
a remainder term
\[
 O\!\left(\frac{x^{1+\varepsilon}}s\right)+O\!\left(x^\varepsilon\right).
\]
Wigert's method, an application of the Euler--Maclaurin sum-formula, involves
rather elaborate calculations, and we have thought it worth while to
indicate an alternative proof, based on the double Fourier formula
\[
 \sum_{p,q=-\infty}^{\infty}f(p,q)
 =\sum_{m,n=-\infty}^{\infty}\int_{-\infty}^{\infty}
  \int_{-\infty}^{\infty}e^{2\pi i(mu+nv)}f(u,v)\dd u\dd v.
\]
This formula is easily seen to be valid when
\[
 f(p,q)=\exp\!\left\{-\left(\frac{p^2+q^2}{x}\right)^s\right\},
\]
and we have
\begin{align*}
&\sum_{p,q=-\infty}^{\infty}
 \exp\!\left\{-\left(\frac{p^2+q^2}{x}\right)^s\right\} \\
&=\sum_{p,q=-\infty}^{\infty}\int_{-\infty}^{\infty}
  \int_{-\infty}^{\infty}
 \cos(2\pi pu)\cos(2\pi qv)
 \exp\!\left\{-\left(\frac{u^2+v^2}{x}\right)^s\right\}\dd u\dd v \\
&=\frac{x}{2}\sum_{p,q=-\infty}^{\infty}\int_0^\infty e^{-R^s}\dd R
 \int_0^{2\pi}
 \cos\!\left(2\pi p\sqrt{xR}\cos\phi\right)
 \cos\!\left(2\pi q\sqrt{xR}\sin\phi\right)\dd\phi \\
&=\pi x\int_0^\infty e^{-R^s}\dd R
 +\pi x\sum_{n=1}^{\infty}r(n)\int_0^\infty e^{-R^s}
 J_0\!\left(2\pi\sqrt{nxR}\right)\dd R\,^{*} \\
&=\pi x\Gamma\!\left(1+\frac1s\right)
 +s\sqrt{x}\sum_{n=1}^{\infty}\frac{r(n)}{\sqrt n}
 \int_0^\infty e^{-R^s}R^{s-1/2}
 J_1\!\left(2\pi\sqrt{nxR}\right)\dd R\,^{\dagger} \\
&=\pi x+\sqrt{x}\sum_{n=1}^{\infty}\frac{r(n)}{\sqrt n}
 \int_0^\infty e^{-u}u^{1/(2s)}
 J_1\!\left(2\pi\sqrt{nx}\,u^{1/(2s)}\right)\dd u
 +O\!\left(x^{37/112}\right).
 \tag{2.2}
\end{align*}
\begin{quote}\small
$^{*}$ See G. N. Watson, \textit{Bessel Functions}, p.~21, \S~2.21,
formula~1.\\
$^{\dagger}$ We integrate by parts, and use Watson, p.~46, \S~3.2,
formula~5.
\end{quote}

On the other hand we have, in virtue of $r(n)=O(n^\varepsilon)$,
\begin{align*}
\mathrm{R}(x)&-\sum_{p,q=-\infty}^{\infty}
 \exp\!\left\{-\left(\frac{p^2+q^2}{x}\right)^s\right\} \\
&=\sum_{0\leq n\leq x}r(n)
 \left\{1-\exp\!\left[-\left(\frac nx\right)^s\right]\right\}
 -\sum_{n>x}r(n)\exp\!\left[-\left(\frac nx\right)^s\right] \\
&=O\!\left(\int_0^x u^\varepsilon\left(\frac ux\right)^s\dd u\right)
 +O\!\left(\int_x^\infty u^\varepsilon
 \exp\!\left[-\left(\frac ux\right)^s\right]\dd u\right)
 +O\!\left(x^\varepsilon\right) \\
&=O\!\left(\frac{x^{1+\varepsilon}}s\right)
 =O\!\left(x^{37/112+\varepsilon}\right).
\end{align*}
Hence, from (2.2),
\begin{equation}
\begin{aligned}
 \mathrm{R}(x)-\pi x
 &=\sqrt{x}\sum_{n=1}^{\infty}\frac{r(n)}{\sqrt n}
 \int_0^\infty e^{-u}u^{1/(2s)}
 J_1\!\left(2\pi\sqrt{nx}\,u^{1/(2s)}\right)\dd u \\
 &\qquad+O\!\left(x^{37/112+\varepsilon}\right).
\end{aligned}
\tag{2.3}
\end{equation}
Since
\[
 y^{3/2}\left|J_1(y)-\sqrt{\frac{2}{\pi y}}
 \cos\!\left(y-\frac{3\pi}{4}\right)\right|
\]
is less than an absolute constant for $y>0$, it follows from (2.3) that
\begin{align*}
\mathrm{R}(x)-\pi x
={}&\frac1\pi x^{1/4}\sum_{n=1}^{\infty}\frac{r(n)}{n^{3/4}}
 \int_0^\infty e^{-u}u^{1/(4s)}
 \cos\!\left(2\pi\sqrt{nx}\,u^{1/(2s)}-\frac{3\pi}{4}\right)\dd u \\
&+O\!\left(x^{-1/4}\right)\sum_{n=1}^{\infty}\frac{r(n)}{n^{5/4}}
 \int_0^\infty e^{-u}u^{-7/(4s)}\dd u
 +O\!\left(x^{37/112+\varepsilon}\right) \\
={}&\frac1\pi x^{1/4}\sum_{n=1}^{\infty}\frac{r(n)}{n^{3/4}}a_n
 +O\!\left(x^{37/112+\varepsilon}\right),
\end{align*}
the desired result.

\medskip
\noindent\textbf{3.} \textsc{Lemma 2.} \textit{We have}
\[
 \mathrm{R}(x)-\pi x
 =\frac1\pi x^{1/4}\sum_{n\leq x^{19/56+\varepsilon}}
 \frac{r(n)}{n^{3/4}}a_n+O\!\left(x^{37/112+\varepsilon}\right).
\]
For each integer $\rho$ and large $x$ (or $s$) we have, writing $u^{2s}$
for $u$ in (2.1) and integrating $2\rho$ times by parts,
\[
 |a_n|=\frac{2s}{(4\pi^2nx)^\rho}
 \left|\int_0^\infty
 \cos\!\left(2\pi\sqrt{nx}\,u-\frac{3\pi}{4}\right)
 \left(\frac{d}{du}\right)^{2\rho}
 \left\{e^{-u^{2s}}u^{2s-1/2}\right\}\dd u\right|.
\]
Integrating once more by parts and observing that
\begin{align*}
&\int_0^\infty\left|
 \left(\frac{d}{du}\right)^{2\rho+1}
 \left\{e^{-u^{2s}}u^{2s-1/2}\right\}\right|\dd u \\
&\qquad<\sum_{m=0}^{2\rho+1}Ks^m\int_0^\infty e^{-u^{2s}}
 \left(u^{2s+Ks}+u^{2s-K}\right)\dd u
 <\sum Ks^m s^{-1}<Ks^{2\rho},
\end{align*}
we conclude that, for a suitable $x(\rho)\geq2$ and $x\geq x(\rho)$,
\[
 |a_n|<K\left(\frac{s^2}{nx}\right)^{\rho+1/2}
 =K\left(\frac{x^{19/56}}n\right)^{\rho+1/2},
\]
where $K$ depends only on $\rho$. From this it follows that
\begin{align*}
&\frac1\pi x^{1/4}\sum_{n>x^{19/56+\varepsilon}}
 \frac{r(n)}{n^{3/4}}|a_n| \\
&\qquad=O\!\left(Kx^{1/4+(19/56)(\rho+1/2)}\right)
 \sum_{n>x^{19/56+\varepsilon}}n^{\varepsilon-3/4-\rho-1/2},
\end{align*}
and this, by a suitable choice of $\rho=\rho(\varepsilon)$ (sufficiently
large), is $O(x^{37/112+\varepsilon})$. Lemma~2 now follows from Lemma~1.

\medskip
\noindent\textbf{4.} \textsc{Lemma 3.} \textit{We have}
\begin{equation}
\begin{aligned}
 \mathrm{R}(x)-\pi x
 &=\frac1\pi x^{1/4}\int_{x^{-1}}^x e^{-u}u^{1/(4s)}
 \sum_{n\leq x^{19/56+\varepsilon}}\frac{r(n)}{n^{3/4}}
 \cos\!\left(z\sqrt n-\frac{3\pi}{4}\right)\dd u \\
 &\qquad+O\!\left(x^{37/112+\varepsilon}\right),
\end{aligned}
\tag{4.1}
\end{equation}
\textit{where}
\[
 z=2\pi\sqrt{x}\,u^{1/(2s)}.
\]
For
\begin{align*}
&\frac1\pi x^{1/4}\sum_{n\leq x^{19/56+\varepsilon}}
 \frac{r(n)}{n^{3/4}}
 \left(\int_0^{x^{-1}}+\int_x^\infty\right)e^{-u}u^{1/(4s)}
 \cos\!\left(2\pi\sqrt{nx}\,u^{1/(2s)}-\frac{3\pi}{4}\right)\dd u \\
&\qquad=O\!\left(x^{1/4}\right)
 \sum_{n\leq x^{19/56+\varepsilon}}
 \left(\int_0^{x^{-1}}\dd u+\int_x^\infty e^{-u}u\dd u\right) \\
&\qquad=O\!\left(x^{1/4}\right)\sum_n x^{-1}
 =O\!\left(x^{37/112+\varepsilon}\right),
\tag{4.2}
\end{align*}
and the result follows from (4.2) and Lemma~2.

\medskip
\noindent\textbf{5.} \textsc{Lemma 4.}
\textit{For $\frac{19}{56}+\varepsilon\leq\frac23$, and uniformly for
$x^{-1}\leq u\leq x$, we have}
\begin{equation}
 \sum_{n\leq x^{19/56+\varepsilon}}\frac{r(n)}{n^{3/4}}e^{iz\sqrt n}
 =O\!\left(x^{9/112+\varepsilon}\right).
 \tag{5.1}
\end{equation}
The main theorem follows at once when Lemma~4 is proved. For then the sum in
(4.1) is $O(x^{9/112+\varepsilon})$, and we have
\begin{align*}
 \mathrm{R}(x)-\pi x
 &=O\!\left(x^{1/4+9/112+\varepsilon}\right)
 \int_{x^{-1}}^x e^{-u}u^{1/(4s)}\dd u \\
 &=O\!\left(x^{37/112+\varepsilon}\right)
 \int_0^\infty e^{-u}(1+u)\dd u
 =O\!\left(x^{37/112+\varepsilon}\right).
\end{align*}

We prove first:

\medskip
\noindent\textsc{Lemma 5.}\footnote{It is here that we employ the method of
Hardy and Littlewood referred to in \S~1.}
\textit{Let $\alpha>0$, $0\leq\tau\leq\frac13$,
\[
 \frac1{\sqrt2}x^\tau\leq Q\leq Q'\leq2x^\tau,
 \qquad 0<p\leq Q,
\]
$Q$ and $Q'$ being integers. Then}
\begin{equation}
 \left|\sum_{q=Q}^{Q'}e^{iz\sqrt{p^2+q^2}}\right|
 \leq c_1\left(x^{9\tau/10+1/80+\alpha}
 +p^{-1/4}x^{49\tau/40-1/80+\alpha}\right).
 \tag{5.2}
\end{equation}
We may suppose that $\tau>\frac18$; for if $\tau\leq\frac18$,
\[
 \left|\sum_{q=Q}^{Q'}e^{iz\sqrt{p^2+q^2}}\right|
 \leq2x^\tau\leq2x^{9\tau/10+1/80+\alpha}.
\]
We write
\[
 \nu=\left[x^{4\tau/5-1/10}\right],\qquad
 N=\left[\frac{Q'-Q}{\nu}\right],\qquad
 v_n=Q+n\nu,
\]
$n$ being a non-negative integer. Clearly $\nu\geq1$. We may suppose $N\geq2$,
since otherwise
\[
 \left|\sum_{q=Q}^{Q'}e^{iz\sqrt{p^2+q^2}}\right|
 \leq\nu\frac{Q'-Q}{\nu}+1\leq3\nu
 \leq3x^{9\tau/10+1/80+\alpha}.
\]
And we may further suppose $x\geq c_2\geq2$, where $c_2$ has the property
\[
 \frac{3\nu}{v_n}\leq\frac{3\nu}{Q}\leq\frac12
 \qquad(x\geq c_2),
\]
evidently valid if $c_2$ is large enough. Since
\[
 v_{N-1}+\nu=v_N>Q+\left(\frac{Q'-Q}{\nu}-1\right)\nu=Q'-\nu,
\]
we have
\begin{equation}
\begin{aligned}
&\left|\sum_{q=Q}^{Q'}e^{iz\sqrt{p^2+q^2}}
 -\sum_{n=0}^{N-1}\sum_{m=0}^{\nu-1}
 e^{iz\sqrt{p^2+(v_n+m)^2}}\right| \\
&\qquad\leq\nu\leq x^{9\tau/10+1/80+\alpha}.
\end{aligned}
\tag{5.3}
\end{equation}
Let now
\[
 \sqrt{p^2+(v_n+y)^2}
 =\sqrt{p^2+v_n^2}\left\{P(y)+y^5f(y)\right\},
\]
where the square root is taken positive for $y\geq0$,
\[
 P(y)=P_n(y)=ay^4+a_1y^3+a_2y^2+a_3y+a_4
\]
is a real polynomial of the fourth degree in $y$, and $f(y)$ is regular at
$y=0$. The number $a$ is given by
\begin{equation}
 a=\frac{p^2(4v_n^2-p^2)}{8(p^2+v_n^2)^4}.
 \tag{5.4}
\end{equation}
We write further
\[
 g(y)=e^{iz\sqrt{p^2+v_n^2}\,y^5f(y)}=\sum_{h=0}^{\infty}g_h y^h,
\]
so that in the neighbourhood of $y=0$
\begin{align*}
 g(y)=\exp\Bigg[iz\sqrt{p^2+v_n^2}\Bigg\{&
 \binom{1/2}{3}\frac{6v_ny^5+y^6}{(p^2+v_n^2)^3} \\
 &+\binom{1/2}{4}
 \frac{32v_n^3y^5+24v_n^2y^6+8v_ny^7+y^8}{(p^2+v_n^2)^4} \\
 &+\sum_{k=5}^{\infty}\binom{1/2}{k}
 \left(\frac{y^2+2v_ny}{p^2+v_n^2}\right)^k\Bigg\}\Bigg].
\end{align*}
Since $z\leq c_3\sqrt{x}$, it follows that
\begin{align*}
 \sum_{h=0}^{\infty}|g_h|\nu^h
 &\leq\exp\left[c_4z\left\{
 \frac{\nu^5v_n}{v_n^5}+\frac{\nu^5v_n^3}{v_n^7}
 +v_n\sum_{k=5}^{\infty}\left(\frac{3\nu v_n}{v_n^2}\right)^k
 \right\}\right] \\
 &\leq\exp\left[c_5\sqrt{x}\frac{\nu^5}{v_n^4}\right]\leq c_6.
\end{align*}
Hence
\begin{align*}
&\left|\sum_{m=0}^{\nu-1}e^{iz\sqrt{p^2+(v_n+m)^2}}\right| \\
&\quad=\left|\sum_{m=0}^{\nu-1}e^{iz\sqrt{p^2+v_n^2}\,P(m)}
 \sum_{h=0}^{\infty}g_hm^h\right| \\
&\quad\leq\sum_{h=0}^{\infty}|g_h|
 \left|\sum_{m=0}^{\nu-1}e^{iz\sqrt{p^2+v_n^2}\,P(m)}m^h\right| \\
&\quad\leq\sum_{h=0}^{\infty}|g_h|\nu^h
 \max_{0\leq\nu'\leq\nu-1}
 \left|\sum_{m=\nu'}^{\nu-1}e^{iz\sqrt{p^2+v_n^2}\,P(m)}\right| \\
&\quad\leq c_6\max_{0\leq\nu'\leq\nu-1}
 \left|\sum_{m=\nu'}^{\nu-1}e^{iz\sqrt{p^2+v_n^2}\,P(m)}\right|.
\tag{5.5}
\end{align*}

\medskip
\noindent\textbf{6.} \textsc{Lemma 6.}
\textit{Suppose $b$ integral and $b\geq2$;
\[
 P(y)=\beta_0y^b+\beta_1y^{b-1}+\cdots+\beta_b
\]
a polynomial with real coefficients; $\psi$ real, $\nu_1$ integral,
$\nu_2$ integral and $\nu_2\geq1$. Then}
\begin{align*}
 \left|\sum_{m=\nu_1+1}^{\nu_1+\nu_2}e^{i\psi P(m)}\right|^{2^{b-1}}
 &\leq(2\nu_2)^{2^{b-1}-b}
 \sum_{|m_1|+\cdots+|m_{b-1}|\leq\nu_2-1} \\
 &\qquad\Min\left\{\nu_2,
 \left|\cosec\!\left(\tfrac12\psi\beta_0b!m_1\cdots m_{b-1}\right)\right|
 \right\}.\footnotemark
\end{align*}
\footnotetext{The ``Min'' denotes $\nu_2$ if the cosecant is infinite.}
This result is due to Weyl.\footnote{H. Weyl, ``Über ein Problem aus dem
Gebiet der Diophantischen Approximationen,'' \textit{Göttinger Nachrichten}
(1914), pp.~234--244 (241--243). See also the same author, ``Über die
Gleichverteilung von Zahlen mod. Eins,'' \textit{Math. Annalen}, vol.~77,
pp.~313--352 (328--331) (1916), and again ``Zur Abschätzung von
$\zeta(1+it)$,'' \textit{Math. Zeitschr.}, vol.~10, pp.~88--101 (90) (1921).
The form we adopt, as sufficient for our purpose, is taken from a paper by
E. Landau, ``Zum Waringschen Problem,'' \textit{Math. Zeitschr.}, vol.~12,
pp.~219--247 (224--225) (1922).}
In Lemma~6 we now take
\[
 b=4,\qquad \nu_1=\nu'-1,\qquad \nu_2=\nu-\nu',\qquad
 \psi=z\sqrt{p^2+v_n^2},\qquad \beta_0=a.
\]
We obtain, writing for brevity
\begin{equation}
 t_n=12az\sqrt{p^2+v_n^2}\,m_1m_2m_3,
 \tag{6.1}
\end{equation}
\begin{align*}
&\left|\sum_{m=\nu'}^{\nu-1}e^{iz\sqrt{p^2+v_n^2}\,P(m)}\right| \\
&\qquad\leq c_7(\nu-\nu')^{1/2}
 \left\{\sum_{|m_1|+|m_2|+|m_3|\leq\nu-\nu'-1}
 \Min\left(\nu-\nu',|\cosec t_n|\right)\right\}^{1/8},
\end{align*}
and consequently
\begin{align*}
&\max_{0\leq\nu'\leq\nu-1}
 \left|\sum_{m=\nu'}^{\nu-1}e^{iz\sqrt{p^2+v_n^2}\,P(m)}\right| \\
&\qquad\leq c_7\nu^{1/2}
 \left\{\sum_{|m_1|+|m_2|+|m_3|\leq\nu}
 \Min\left(\nu,|\cosec t_n|\right)\right\}^{1/8} \\
&\qquad\leq c_8\nu^{1/2}
 \left\{\nu^3+\sum_{m_1,m_2,m_3=1}^{\nu}
 \Min\left(\nu,|\cosec t_n|\right)\right\}^{1/8} \\
&\qquad\leq c_8\left[\nu^{7/8}+\nu^{1/2}
 \left\{\sum_{m_1,m_2,m_3=1}^{\nu}
 \Min\left(\nu,|\cosec t_n|\right)\right\}^{1/8}\right].
\end{align*}
From this and (5.5) it follows that
\begin{align*}
&\left|\sum_{n=0}^{N-1}\sum_{m=0}^{\nu-1}
 e^{iz\sqrt{p^2+(v_n+m)^2}}\right| \\
&\qquad\leq c_9\left[N\nu^{7/8}+\nu^{1/2}\sum_{n=0}^{N-1}
 \left\{\sum_{m_1,m_2,m_3=1}^{\nu}
 \Min\left(\nu,|\cosec t_n|\right)\right\}^{1/8}\right] \\
&\qquad\leq c_9\left[N\nu^{7/8}+N^{7/8}\nu^{1/2}
 \left\{\sum_{m_1,m_2,m_3=1}^{\nu}\sum_{n=0}^{N-1}
 \Min\left(\nu,|\cosec t_n|\right)\right\}^{1/8}\right],
\tag{6.2}
\end{align*}
by the generalised Schwarz inequality. By the hypotheses of Lemma~5,
\[
 N\nu^{7/8}\leq c_{10}x^{9\tau/10+1/80},\qquad
 N^{7/8}\nu^{1/2}\leq c_{10}x^{23\tau/40+3/80}.
\]
Hence it is easily verified, from (5.3) and (6.2), that Lemma~5 follows when
we have shown that
\begin{equation}
 \sum_{m_1,m_2,m_3=1}^{\nu}\sum_{n=0}^{N-1}
 \Min\left(\nu,|\cosec t_n|\right)
 \leq c_{11}\left(x^{13\tau/5-1/5+\alpha}
 +p^{-2}x^{26\tau/5-2/5+2\alpha}\right).
 \tag{6.3}
\end{equation}
Now from (5.4) and (6.1) we have
\[
 t_n=\frac32zm_1m_2m_3p^2
 \frac{4v_n^2-p^2}{(p^2+v_n^2)^{7/2}}.
\]
Since
\[
 \frac{d}{dy}\frac{4y^2-p^2}{(p^2+y^2)^{7/2}}
 =\frac{5y(3p^2-4y^2)}{(p^2+y^2)^{9/2}},
\]
and
\[
 3p^2\leq3Q^2\leq3v_n^2,\qquad
 v_{N-1}<v_N\leq Q+\frac{Q'-Q}{\nu}\nu=Q'\leq2x^\tau,
 \qquad z>c_{12}\sqrt{x},
\]
we have for $0\leq n\leq N-2$
\begin{align*}
 t_n-t_{n+1}
 &\geq(v_{n+1}-v_n)\frac32zm_1m_2m_3p^2
 \min_{v_n\leq y\leq v_{n+1}}
 \frac{5y(4y^2-3p^2)}{(p^2+y^2)^{9/2}} \\
 &\geq c_{13}\nu\sqrt{x}\,m_1m_2m_3p^2
 \frac{v_n^3}{v_{n+1}^9} \\
 &\geq c_{14}m_1m_2m_3p^2x^{2/5-26\tau/5}.
\tag{6.4}
\end{align*}
Let now $\xi_g$ ($g$ integral) denote the interval
\[
 \frac12\pi g<y\leq\frac12\pi(g+1).
\]
We consider the part sum
\[
 \eta_g=\sum_g\Min\left(\nu,|\cosec t_n|\right)=\sum_g u_n,
\]
where $\sum_g$ is taken over those values of $n$ between $0$ and $N-1$
inclusive for which $t_n$ lies in $\xi_g$. Let $g^*$ be $\frac12g$ or
$\frac12(g+1)$, whichever is an integer, and $n^*$ the rank of that term of
$\eta_g$ for which $|t_{n^*}-\pi g^*|$ is a minimum. For $n=n^*$ we have
$u_n\leq\nu$; for other values of $n$, $u_n\leq|\cosec t_n|$, and it follows
from (6.4) that
\begin{equation}
 \eta_g\leq\nu+\sum_{\substack{t_n\in\xi_g\\n\neq n^*}}|\cosec t_n|
 \leq\nu+\sum_k'\left|\cosec\!\left(
 c_{14}m_1m_2m_3p^2x^{2/5-26\tau/5}k\right)\right|,
 \tag{6.5}
\end{equation}
where $\sum'$ is taken only over $k$ for which
\[
 1\leq k\leq N,\qquad
 c_{14}m_1m_2m_3p^2x^{2/5-26\tau/5}k\leq\frac12\pi.
\]
From this it follows that
\begin{equation}
 \eta_g\leq\nu+c_{15}\frac{x^{26\tau/5-2/5}}{m_1m_2m_3p^2}
 \sum_{k=1}^{N}\frac1k
 \leq c_{16}\left(\nu+
 \frac{x^{26\tau/5-2/5+\alpha}}{m_1m_2m_3p^2}\right).
 \tag{6.6}
\end{equation}
Again, from (5.4) and (6.1),
\[
 t_n\leq c_{17}m_1m_2m_3p^2x^{1/2-5\tau}.
\]
The number of intervals $\xi_g$ (concerned with $t_n$'s) therefore does not
exceed
\[
 c_{18}\left(m_1m_2m_3p^2x^{1/2-5\tau}+1\right).\footnotemark
\]
\footnotetext{The term $1$ cannot be omitted, since
$x^{1/2-5\tau}\to0$ as $x\to\infty$; it has a real importance.}
Hence, from (6.6),
\begin{align*}
&\sum_{m_1,m_2,m_3=1}^{\nu}\sum_{n=0}^{N-1}
 \Min\left(\nu,|\cosec t_n|\right) \\
&\quad\leq c_{19}\sum_{m_1,m_2,m_3=1}^{\nu}
 \left(m_1m_2m_3p^2x^{1/2-5\tau}+1\right)
 \left(\nu+\frac{x^{26\tau/5-2/5+\alpha}}{m_1m_2m_3p^2}\right) \\
&\quad\leq c_{19}\left\{p^2\nu^7x^{1/2-5\tau}+\nu^4
 +\nu^3x^{\tau/5+1/10+\alpha}
 +p^{-2}x^{26\tau/5-2/5+\alpha}
 \left(\sum_{m=1}^{\nu}\frac1m\right)^3\right\} \\
&\quad\leq c_{20}\left(x^{13\tau/5-1/5}+x^{16\tau/5-2/5}
 +x^{13\tau/5-1/5+\alpha}
 +p^{-2}x^{26\tau/5-2/5+2\alpha}\right) \\
&\quad\leq c_{21}\left(x^{13\tau/5-1/5+\alpha}
 +p^{-2}x^{26\tau/5-2/5+2\alpha}\right).
\end{align*}
This is (6.3), and Lemma~5 is established.

\medskip
\noindent\textbf{7.}
Returning now to Lemma~4, let $0\leq\tau\leq\frac13$. We take
\[
\begin{array}{ll}
 Q=\left[\sqrt{x^{2\tau}-p^2}\right]+1,
 &\displaystyle\left(0<p\leq\frac1{\sqrt2}x^\tau\right),\\[6pt]
 Q=p+1,
 &\displaystyle\left(\frac1{\sqrt2}x^\tau<p\leq\sqrt2x^\tau\right),
\end{array}
\]
and (i.e. in either case) we take
\[
 Q'=\left[\sqrt{4x^{2\tau}-p^2}\right]
 \qquad(0<p\leq\sqrt2x^\tau).
\]
Then
\begin{align*}
&\left|\sum_{x^{2\tau}<n\leq4x^{2\tau}}r(n)e^{iz\sqrt n}
 -8\sum_{0<p\leq\sqrt2x^\tau}\sum_{Q\leq q\leq Q'}
 e^{iz\sqrt{p^2+q^2}}\right| \\
&\qquad=4\left|\sum_{x^\tau<q\leq2x^\tau}e^{izq}
 +\sum_{x^\tau/\sqrt2<q\leq\sqrt2x^\tau}e^{izq\sqrt2}\right|
 \leq c_{22}x^\tau.\footnotemark
\tag{7.1}
\end{align*}
\footnotetext{Every $q$-sum for which $Q>Q'$ is, of course, interpreted to
be zero.}
Now to every existent $q$-sum in the left side of (7.1) corresponds a set of
numbers $p,Q,Q'$ satisfying the conditions of Lemma~5, which gives therefore
\begin{align*}
&\left|\sum_{0<p\leq\sqrt2x^\tau}\sum_{Q\leq q\leq Q'}
 e^{iz\sqrt{p^2+q^2}}\right| \\
&\qquad\leq c_1\sum_{0<p\leq\sqrt2x^\tau}
 \left(x^{9\tau/10+1/80+\alpha}
 +p^{-1/4}x^{49\tau/40-1/80+\alpha}\right) \\
&\qquad\leq c_{23}\left(x^{19\tau/10+1/80+\alpha}
 +x^{79\tau/40-1/80+\alpha}\right)
 \leq c_{24}x^{19\tau/10+1/80+\alpha}.
\end{align*}
Hence, from (7.1),
\begin{equation}
 \left|\sum_{x^{2\tau}<n\leq4x^{2\tau}}r(n)e^{iz\sqrt n}\right|
 \leq c_{25}x^{19\tau/10+1/80+\alpha}.
 \tag{7.2}
\end{equation}
We write now
\[
 x^\tau=x^{\tau_0}=2x^{\tau_1}=4x^{\tau_2}=\cdots=2^M x^{\tau_M},
\]
where $M$ is determined by
\[
 2>x^{\tau_M}\geq1.
\]
Taking $\tau=\tau_m$ in (7.2), we have for $M\geq1$
\begin{align*}
&\left|\sum_{m=1}^{M}\sum_{x^{2\tau_m}<n\leq4x^{2\tau_m}}
 r(n)e^{iz\sqrt n}\right| \\
&\qquad\leq c_{25}\sum_{m=1}^{M}x^{19\tau_m/10+1/80+\alpha}
 \leq c_{25}Mx^{19\tau/10+1/80+\alpha}
 \leq c_{26}x^{19\tau/10+1/80+2\alpha},
\end{align*}
or
\[
 \left|\sum_{1\leq n\leq x^{2\tau}}r(n)e^{iz\sqrt n}\right|
 \leq c_{26}x^{19\tau/10+1/80+2\alpha}+c_{27}
 \leq c_{28}x^{19\tau/10+1/80+2\alpha}.
\]
Writing $y$ for $x^{2\tau}$ in this, we have
\[
 \left|\sum_{1\leq n\leq y}r(n)e^{iz\sqrt n}\right|
 \leq c_{28}x^{1/80+2\alpha}y^{19/20}
 \qquad(1\leq y\leq x^{2/3}).
\]
It follows by partial summation that
\begin{align*}
&\left|\sum_{1\leq n\leq x^{19/56+\varepsilon}}
 \frac{r(n)}{n^{3/4}}e^{iz\sqrt n}\right| \\
&\qquad\leq c_{29}x^{1/80+2\alpha}
 \sum_{1\leq n\leq x^{19/56+\varepsilon}}n^{19/20-7/4}
 \leq c_{30}x^{9/112+\varepsilon/5+2\alpha}.
\end{align*}
If in this we take $\alpha=\frac25\varepsilon$, we obtain the result of
Lemma~4, and our proof is completed.

\bigskip
\begin{center}
\textit{Note on the Preceding Paper by Prof. E. Landau.---Added August 2,
1924.}
\end{center}

\noindent\textbf{8.}
The authors deduce, by means of the ideas of Lemma~5, their main result from
the formula of Lemma~1. It is possible to start from the relation
\begin{equation}
 \int_0^x\{\mathrm{R}(y)-\pi y\}\dd y
 =-c_{31}x^{3/4}\sum_{n=1}^{\infty}\frac{r(n)}{n^{5/4}}
 \cos\!\left(2\pi\sqrt{nx}-\frac\pi4\right)+O\!\left(x^{1/4}\right),
 \tag{8.1}
\end{equation}
which is classical. I am able, indeed, to refine on the authors'
$x^\varepsilon$, and in fact to deduce from (8.1) that
\[
 \mathrm{R}(x)-\pi x
 =O\!\left(x^{37/112}\log^{5/56}x\right).
\]
The proof of Lemma~5 provides the inequality
\[
 \left|\sum_{q=Q}^{Q'}e^{2\pi i\sqrt y\sqrt{p^2+q^2}}\right|
 <c_{32}\left(y^{1/80}Q^{9/10}\log^{1/8}Q
 +y^{-1/80}p^{-1/4}Q^{49/40}\log^{1/2}Q\right)
\]
for $y>1$, $1\leq p<Q\leq Q'<2Q$, whence, by Abel's Lemma, with
$j=\frac34$ or $\frac54$,
\begin{equation}
\begin{aligned}
&\left|\sum_{q=Q}^{Q'}
 \frac{e^{2\pi i\sqrt y\sqrt{p^2+q^2}}}{(p^2+q^2)^j}\right| \\
&\quad<c_{32}\left(y^{1/80}Q^{9/10-2j}\log^{1/8}Q
 +y^{-1/80}p^{-1/4}Q^{49/40-2j}\log^{1/2}Q\right).
\end{aligned}
\tag{8.2}
\end{equation}
Now (8.2) is valid without the assumption $Q'<2Q$, provided we replace
$c_{32}$ by $c_{33}$.

Therefore, if $2\leq\omega\leq y^{1/2}$,
\begin{align*}
&\left|\sum_{n\leq\omega}\frac{r(n)}{n^{3/4}}e^{2\pi i\sqrt{ny}}\right| \\
&\quad<c_{34}+8\sum_{1<p\leq\sqrt{\omega/2}}
 \left|\sum_{p<q\leq\sqrt{\omega-p^2}}
 \frac{e^{2\pi i\sqrt y\sqrt{p^2+q^2}}}{(p^2+q^2)^{3/4}}\right| \\
&\quad\leq c_{34}+c_{35}\sum_{1\leq p\leq\sqrt{\omega/2}}
 \left(y^{1/80}p^{-3/5}\log^{1/8}p
 +y^{-1/80}p^{-21/40}\log^{1/2}p\right) \\
&\quad<c_{36}\left(y^{1/80}\omega^{1/5}\log^{1/8}y
 +y^{-1/80}\omega^{19/80}\log^{1/2}y\right) \\
&\quad<c_{37}y^{1/80}\omega^{1/5}\log^{1/8}y,
\tag{8.3}
\end{align*}
and
\begin{align*}
&\left|\sum_{n>\omega}\frac{r(n)}{n^{5/4}}e^{2\pi i\sqrt{ny}}\right| \\
&\quad<c_{38}\omega^{-3/4}+8\sum_{p=1}^{\infty}
 \left|\sum_{q>\Max\{\sqrt{\omega-p^2},p\}}
 \frac{e^{2\pi i\sqrt y\sqrt{p^2+q^2}}}{(p^2+q^2)^{5/4}}\right| \\
&\quad<c_{38}\omega^{-3/4}
 +c_{39}\sum_{1\leq p\leq\sqrt{\omega/2}}
 \left(y^{1/80}\omega^{-4/5}\log^{1/8}\omega
 +y^{-1/80}p^{-1/4}\omega^{-51/80}\log^{1/2}\omega\right) \\
&\qquad+c_{40}\sum_{p>\sqrt{\omega/2}}
 \left(y^{1/80}p^{-8/5}\log^{1/8}p
 +y^{-1/80}p^{-61/40}\log^{1/2}p\right) \\
&\quad<c_{41}\left(\omega^{-3/4}
 +y^{1/80}\omega^{-3/10}\log^{1/8}y
 +y^{-1/80}\omega^{-21/80}\log^{1/2}y\right) \\
&\quad<c_{42}y^{1/80}\omega^{-3/10}\log^{1/8}y.
\tag{8.4}
\end{align*}
I now take
\[
 \omega=x^{19/56}\log^{-5/28}x,\qquad
 z=x^{37/112}\log^{5/56}x.
\]
Observing that
\begin{align*}
&\int_x^{x\pm z}\sum_{n\leq\omega}\frac{r(n)}{n^{5/4}}
 \frac{d}{dy}\left\{y^{3/4}
 \cos\!\left(2\pi\sqrt{ny}-\frac\pi4\right)\right\}\dd y \\
&\quad=-c_{43}\int_x^{x\pm z}y^{1/4}
 \sum_{n\leq\omega}\frac{r(n)}{n^{3/4}}
 \sin\!\left(2\pi\sqrt{ny}-\frac\pi4\right)\dd y
 +O\!\left(zx^{-1/4}\right),
\end{align*}
we obtain, from (8.1), (8.3) and (8.4),
\begin{align*}
 \int_x^{x\pm z}\{\mathrm{R}(y)-\pi y\}\dd y
 &=\int_x^{x\pm z}\mathrm{R}(y)\dd y\mp\pi xz+O(z^2) \\
 &=O\!\left(zx^{1/4}x^{1/80}\omega^{1/5}\log^{1/8}x\right)
   +O\!\left(zx^{-1/4}\right) \\
 &\quad+O\!\left(x^{3/4}x^{1/80}\omega^{-3/10}\log^{1/8}x\right)
   +O\!\left(x^{1/4}\right) \\
 &=O(z^2),
\end{align*}
and hence
\[
 \pi x+O(z)\leq\frac1z\int_{x-z}^{x}\mathrm{R}(y)\dd y
 \leq\mathrm{R}(x)\leq\frac1z\int_x^{x+z}\mathrm{R}(y)\dd y
 \leq\pi x+O(z).
\]

\end{document}
